\documentclass[10pt]{article}
\usepackage[T1]{fontenc}
\usepackage[utf8]{inputenc}
\usepackage[french]{babel}
\usepackage{amsmath,amssymb,mathrsfs,mathpazo}
\usepackage{geometry}
\geometry{a4paper,margin=2cm}
\usepackage{graphicx,array,tabularx,tabularray,colortbl,multirow,multicol,booktabs,enumitem,pifont,ulem,fancybox,diagbox,cancel}
\usepackage{pstricks,pstricks-add,pst-plot,pst-tree,pst-node,pst-text,pst-func,pst-math,pst-bspline,pst-3dplot,pst-eucl,pst-solides3d}
\let\pscalculate\undefined
\usepackage{pst-bezier}
\usepackage{tikz,pgfplots,tkz-tab}
\usetikzlibrary{arrows,arrows.meta,patterns,calc,positioning,trees,decorations.markings,shapes.geometric,shapes.misc,shapes.arrows,intersections,decorations.pathmorphing,decorations.pathreplacing,backgrounds,shadows,babel}
\pgfplotsset{compat=1.18}
\usepgfplotslibrary{fillbetween}
\usepackage[np]{numprint}
\usepackage{listings,makecell,pas-tableur,twoopt}
\let\barre\undefined
\usepackage[pstricks]{tablvar}
\definecolor{darkgreen}{RGB}{0,100,0}
\providecommand{\red}{\color{red}}
\providecommand{\blue}{\color{blue}}
\providecommand{\green}{\color{green}}
\providecommand{\black}{\color{black}}
\providecommand{\white}{\color{white}}
\providecommand{\gray}{\color{gray}}
\providecommand{\cyan}{\color{cyan}}
\providecommand{\magenta}{\color{magenta}}
\providecommand{\yellow}{\color{yellow}}
\providecommand{\violet}{\color{violet}}
\providecommand{\orange}{\color{orange}}
\providecommand{\euro}{\texteuro}
\providecommand{\EUR}{\texteuro}
\providecommand{\ieme}{\textsuperscript{e}}
\providecommand{\degres}{\ensuremath{^\circ}}
\providecommand{\degree}{\ensuremath{^\circ}}
\providecommand{\vv}[1]{\overrightarrow{#1}}
\providecommand{\Coord}[1]{#1}
\providecommand{\decosix}{}
\providecommand{\starredbullet}{\ensuremath{\bullet}}
\pagestyle{empty}
\setlength{\parindent}{0pt}
\usepackage{adjustbox}\DeclareUnicodeCharacter{2236}{:}\DeclareUnicodeCharacter{2190}{\ensuremath{\leftarrow}}
\providecommand{\R}{\mathbb{R}}
\providecommand{\N}{\mathbb{N}}
\providecommand{\Z}{\mathbb{Z}}
\providecommand{\Q}{\mathbb{Q}}
\providecommand{\D}{\mathbb{D}}
\providecommand{\C}{\mathbb{C}}
\providecommand{\e}{{\mathrm{e}}}
\providecommand{\vect}[1]{\overrightarrow{#1}}
\providecommand{\vectt}[1]{\overrightarrow{\mathrm{#1}}}
\providecommand{\vvt}[1]{\overrightarrow{\mathrm{#1}}}
\providecommand{\np}[2][]{{#2~#1}}
\providecommand{\vv}[1]{{\overrightarrow{#1}}}
\renewcommand{\d}{{\mathrm{d}}}
\providecommand{\degres}{{^\circ}}
\providecommand{\degre}{{^\circ}}
\providecommand{\blue}{\color{blue}}
\providecommand{\red}{\color{red}}
\providecommand{\enskip}{\quad}
\providecommand{\llbracket}{\lbrack\!\lbrack}
\providecommand{\rrbracket}{\rbrack\!\rbrack}
\providecommand{\euro}{{\text{€}}}
\providecommand{\no}{{\text{n°}}}
\providecommand{\interval}[2]{{[#1;#2]}}
\providecommand{\intervFO}[2]{{[#1;#2[}}
\providecommand{\dotfill}{\ldots}

\providecommand{\ds}{\displaystyle}
\providecommand{\textding}[1]{#1}
\providecommand{\Oij}{\(\mathrm{O};\vec{\imath},\vec{\jmath}\)}
\providecommand{\Oijk}{\(\mathrm{O};\vec{\imath},\vec{\jmath},\vec{k}\)}
\providecommand{\Ouv}{\(\mathrm{O};\vec{u},\vec{v}\)}
\providecommand{\intervFF}[2]{\ensuremath{[#1;#2]}}
\renewcommand{\euro}{\ifmmode\text{\texteuro}\else\texteuro\fi}
\DeclareUnicodeCharacter{00B0}{\ifmmode^{\circ}\else\textdegree\fi}
\DeclareUnicodeCharacter{20AC}{\ifmmode\text{\texteuro}\else\texteuro\fi}
% Énoncés : Éducation nationale. Transcription et corrigés : APMEP, Denis Vergès et contributeurs. Découpage et classement : CoopMaths, Sébastien Lozano, Éric Elter et contributeurs. Adaptation de la présentation : Maths Explorer.
% CC BY-SA : https://coopmaths.fr/www/about/
% https://coopmaths.fr/alea/static/bac/2021/tex/bac_2021_06_asie_4.tex
\begin{document}
\noindent\textit{APMEP / CoopMaths. CC BY-SA. Présentation adaptée par Maths Explorer.}\par\medskip
\medskip

\begin{tabular}{|l|}\hline
\textbf{Principaux domaines abordés}\\
-- convexité\\
-- fonction logarithme\\ \hline
\end{tabular}

\bigskip

\textbf{Partie I : lectures graphiques}

\medskip

$f$ désigne une fonction définie et dérivable sur $\R$.

On donne ci-dessous la courbe représentative de la fonction dérivée $f'$.

\begin{center}
\psset{xunit=0.9cm,yunit=3cm}
\begin{pspicture}(-7,-0.8)(7,1.2)
\multido{\n=-7+1}{15}{\psline[linestyle=dashed,linewidth=0.1pt](\n,-0.8)(\n,1.2)}
\multido{\n=-0.8+0.2}{11}{\psline[linestyle=dashed,linewidth=0.1pt](-7,\n)(7,\n)}
\psaxes[linewidth=1.25pt,labelFontSize=\scriptstyle]{->}(0,0)(-7,-0.8)(7,1.2)
\psplot[plotpoints=2000,linewidth=1.25pt,linecolor=red]{-7}{7}{x 2 mul 1 add x dup mul x add 2.5 add div}
\uput[u](0,1){\red Courbe de la fonction dérivée $f'$}
\end{pspicture}
\end{center}

\medskip

\emph{Avec la précision permise par le graphique, répondre aux questions suivantes}

\medskip

\begin{enumerate}
\item Déterminer le coefficient directeur de la tangente à la courbe de la fonction $f$ en $0$.
\item 
	\begin{enumerate}
		\item Donner les variations de la fonction dérivée $f'$.
		\item En déduire un intervalle sur lequel $f$ est convexe.
	\end{enumerate}
\end{enumerate}

\bigskip

\textbf{Partie II : étude de fonction}

\medskip

La fonction $f$ est définie sur $\R$ par 

\[f(x) = \ln \left(x^2 + x + \dfrac{5}{2}\right).\]

\medskip

\begin{enumerate}
\item Calculer les limites de la fonction $f$ en $+\infty$ et en $-\infty$.
\item Déterminer une expression $f'(x)$ de la fonction dérivée de $f$ pour tout $x \in \R$.
\item En déduire le tableau des variations de $f$. On veillera à placer les limites dans ce tableau.
\item 
	\begin{enumerate}
		\item Justifier que l'équation $f(x) = 2$ a une unique solution $\alpha$ dans l'intervalle $\left[-\dfrac{1}{2}~;~+ \infty\right[$. 
		\item Donner une valeur approchée de $\alpha$ à $10^{-1}$ près.
 	\end{enumerate}
\item La fonction $f'$ est dérivable sur $\R$. On admet que, pour tout $x \in  \R$,\,  $f''(x) = \dfrac{-2x^2 - 2x + 4}{\left(x^2 + x + \dfrac{5}{2}\right)^2}$.

Déterminer le nombre de points d'inflexion de la courbe représentative de $f$.
\end{enumerate}

\bigskip

\end{document}
